\documentclass[tikz,border=6pt]{standalone}
% 공통 프리앰블 — PM Day1 교재 그림 (TikZ standalone)
\usepackage{fontspec}
\setmainfont{Apple SD Gothic Neo}
\setsansfont{Apple SD Gothic Neo}
\setmonofont{Menlo}
\usepackage{amssymb}
\usepackage{tikz}
\usetikzlibrary{arrows.meta,positioning,calc,shapes.geometric,fit,backgrounds,matrix,decorations.pathreplacing}
\definecolor{navy}{HTML}{1F3A5F}
\definecolor{teal}{HTML}{2A7F62}
\definecolor{rust}{HTML}{C8553D}
\definecolor{sand}{HTML}{E8E1D5}
\definecolor{ink}{HTML}{222222}
\definecolor{grey}{HTML}{7A7A7A}
\definecolor{lightgrey}{HTML}{EFEFEF}
\definecolor{navylight}{HTML}{DCE6F2}
\definecolor{teallight}{HTML}{DDF0E8}
\definecolor{rustlight}{HTML}{F6E0DA}
\tikzset{
  every node/.style={font=\small, text=ink},
  box/.style={draw=navy, thick, rounded corners=2pt, align=center, inner sep=5pt, fill=white},
  boxfill/.style={box, fill=navylight},
  boxteal/.style={draw=teal, thick, rounded corners=2pt, align=center, inner sep=5pt, fill=teallight},
  boxrust/.style={draw=rust, thick, rounded corners=2pt, align=center, inner sep=5pt, fill=rustlight},
  boxgrey/.style={draw=grey, thick, rounded corners=2pt, align=center, inner sep=5pt, fill=lightgrey},
  arr/.style={-{Stealth[length=2.5mm]}, thick, navy},
  arrteal/.style={-{Stealth[length=2.5mm]}, thick, teal},
  arrrust/.style={-{Stealth[length=2.5mm]}, thick, rust},
  arrgrey/.style={-{Stealth[length=2.5mm]}, thick, grey},
  lbl/.style={font=\footnotesize, text=grey, align=center},
  src/.style={font=\scriptsize, text=grey, align=left},
  title/.style={font=\normalsize\bfseries, text=navy, align=center},
}

\begin{document}
\begin{tikzpicture}[node distance=3mm]
\node[font=\small\bfseries, text=navy] (h6) at (0,6mm) {6판 (2017) — 10 Knowledge Areas};
\node[font=\small\bfseries, text=teal] (h8) at (80mm,6mm) {8판 (2025) — 7 Performance Domains};
\foreach \i/\t in {1/Integration,2/Scope,3/Schedule,4/Cost,5/Quality,6/Resource,7/Communications,8/Risk,9/Procurement,10/Stakeholder}{
  \node[boxfill, text width=34mm, minimum height=7mm, font=\footnotesize, anchor=north] (k\i) at (0,{-(\i-1)*9.5mm}) {\t};
}
\foreach \i/\t in {1/Governance,2/Scope,3/Schedule,4/Finance,5/Stakeholders,6/Resources,7/Risk}{
  \node[boxteal, text width=34mm, minimum height=7mm, font=\footnotesize, anchor=north] (d\i) at (80mm,{-(\i-1)*13mm}) {\t};
}
% 매핑
\draw[arr] (k1) -- (d1);
\draw[arr] (k2) -- (d2);
\draw[arr] (k3) -- (d3);
\draw[arr] (k4) -- (d4);
\draw[arrrust, dashed] (k5.east) -- (d1.west);
\draw[arrrust, dashed] (k5.east) -- (d2.west);
\draw[arr] (k6) -- (d6);
\draw[arrrust, dashed] (k7.east) -- (d5.west);
\draw[arr] (k8) -- (d7);
\draw[arrrust, dashed] (k9.east) -- ++(12mm,0) node[boxrust, anchor=west, font=\scriptsize, text width=26mm] {독립 도메인 아님 —\\배치는 인쇄본으로 확인};
\draw[arr] (k10) -- (d5);
\node[lbl, font=\scriptsize, anchor=west, align=left] at (d1.east) {\ Integration 개칭·승격 — 결정 권한의 구조\\\ + Quality의 QA(프로세스 감사)};
\node[lbl, font=\scriptsize, anchor=west, align=left] at (d2.east) {\ + Quality의 QC(산출물 검사)·인수};
\node[lbl, font=\scriptsize, anchor=west, align=left] at (d4.east) {\ Cost 개칭·확장 — 자금 조달·현금흐름};
\node[lbl, font=\scriptsize, anchor=west, align=left] at (d5.east) {\ Stakeholder + Communications 흡수};
\node[lbl, font=\scriptsize, anchor=west, align=left] at (d6.east) {\ 7판 Team에서 이름 회귀};
\node[lbl, font=\scriptsize, anchor=west, align=left] at (d7.east) {\ 7판 Uncertainty에서 이름 회귀};
% 범례
\node[lbl, anchor=north west, align=left, text width=60mm] at (100mm,-86mm) {실선: 조항 번호만 교체하면 되는 매핑\\파선(붉음): 절차 본문의 재배치가 필요한 실질 변경\\[3pt]사내 표준 재매핑 대장의 4열:\\사내 문서·조항 / 6판 근거 / 8판 대응 위치 / 조치};
\node[src, anchor=north west] at (-20mm,-104mm) {출처: PMBOK Guide 6판(PMI, 2017)과 8판(PMI, 2025)의 구조를 교재 2.1\textasciitilde 2.2절의 서술에 따라 재구성. Quality·Communications·Procurement의 정확한 배치는 8판 인쇄본 대조 필요.};
\end{tikzpicture}
\end{document}
